\begin{table}[!ht]
\centering
\caption{Goods-to-services: contemporaneous and lagged connectedness (quarterly)}
\label{tab:h5q}
\setlength{\tabcolsep}{6pt}\footnotesize
\begin{threeparttable}
\begin{tabular}{c c c c c}
\toprule
 & \multicolumn{2}{c}{Goods $\to$ Services} & \multicolumn{2}{c}{Services $\to$ Goods} \\
\cmidrule(lr){2-3}\cmidrule(lr){4-5}
$\tau$ & Contemp. & Lagged & Contemp. & Lagged \\
\midrule
0.5 & 3.3 & 9.9 & 3.4 & 4.4 \\
0.7 & 8.5 & 11.2 & 8.7 & 7.1 \\
0.9 & 13.1 & 18.4 & 15.2 & 18.3 \\
\bottomrule
\end{tabular}
\begin{tablenotes}[flushleft]\footnotesize
\item \textit{Notes:} Pairwise pseudo-quantile $R^2$ connectedness (\%) between core goods and core services (shelter $+$ core services ex-shelter) on the ex-Great-Inflation quarterly sample (1983Q1--2026Q2, 174 quarters), $n_{\text{lag}}=2$ (a six-month lag horizon). ``Goods $\to$ Services'' is the share of services' variation explained by goods, split into contemporaneous and lagged; ``Services $\to$ Goods'' is the reverse. The \emph{lagged} Goods $\to$ Services term is the ``goods lead'' propagation channel; it exceeds the lagged Services $\to$ Goods term at every quantile (by 4.0 at $\tau=0.7$), consistent with goods leading services. Six months is the horizon the monthly two-lag metric could not span. Dependence, not identified causation.
\end{tablenotes}
\end{threeparttable}
\end{table}

